\subsection{n 阶线性方程}

所谓 n 阶线性微分方程，是指形如

\[
\mathrm{D} ^ {n} x - a _ {1} (t) \mathrm{D} ^ {n - 1} x - \dots - a _ {n - 1} (t) \mathrm{D} x - a _ {n} (t) x = b (t)\tag{33}
\]

的方程，其中 \(a_{k}\) ( \(1 \leqslant k \leqslant n\) ) 与 b 是实变量 t 的实（复）值函数，定义在 R 的一个区间 J 上。IV, p.~164 的一般方法表明，这个方程等价于由 n 个一阶方程组成的线性方程组

\[
\left\{ \begin{array}{l l} \frac {d x _ {k}}{d t} = x _ {k + 1} & (1 \leqslant k \leqslant n - 1) \\ \frac {d x _ {n}}{d t} = a _ {1} (t) x _ {n} + a _ {2} (t) x _ {n - 1} + \dots + a _ {n} (t) x _ {1} + b (t) \end{array} \right.\tag{34}
\]

也就是说，等价于线性方程

\[
\frac {d \mathbf {x}}{d t} = A (t). \mathbf {x} + \mathbf {b} (t)\tag{35}
\]

其中我们令 \(\mathbf{x}=(x_{1},x_{2},\ldots,x_{n})\in\mathbf{C}^{n},\mathbf{b}(t)=(0,0,\ldots,0,b(t))\) ，而矩阵 \(A(t)\) 由下式定义

\[
A (t) = \left( \begin{array}{c c c c c} 0 & 1 & 0 & \ldots & 0 \\ 0 & 0 & 1 & \ldots & 0 \\ \vdots & \vdots & \vdots & \ldots & \vdots \\ 0 & 0 & 0 & \ldots & 1 \\ a _ {n} (t) & a _ {n - 1} (t) & a _ {n - 2} (t) & \ldots & a _ {1} (t) \end{array} \right).
\]

因此，对 n 阶线性方程的研究就是把上面的一般结果应用于特殊的线性方程 (35)。对使函数 \(a_{j}\) ( \(1 \leqslant j \leqslant n\) ) 与 b 均为规则函数的每个区间 J，存在唯一的函数 u，它定义在 J 上，具有 n - 1 阶连续导数，并且在这个区间上（一个可数集的点除外）具有 n 阶规则导数，在 J 的一个可数子集的余集上满足 (33)，并且使得

\[
u (t _ {0}) = x _ {0}, \qquad \mathrm{D} u (t _ {0}) = x _ {0} ^ {\prime}, \dots , D ^ {n - 1} u (t _ {0}) = x _ {0} ^ {(n - 1)}\tag{36}
\]

其中 \(t_0\) 是 \(J\) 的任意一点，而 \(x_0, x_0', \ldots, x_0^{(n-1)}\) 是 \(n\) 个任意的复数。

为使齐次方程的 \(p\) 个积分 \(u_{j}\) ( \(1 \leqslant j \leqslant p\) )

\[
\mathrm{D} ^ {n} x - a _ {1} (t) \mathrm{D} ^ {n - 1} x - \dots - a _ {n - 1} (t) \mathrm{D} x - a _ {n} (t) x = 0\tag{37}
\]

（这个方程与 (33) 相伴）线性无关（在由 J 到 C 的连续映射组成、视为 C 上向量空间的空间 \(\mathcal{C}(\mathrm{J};\mathbf{C})\) 中），必须且只需相应的 p 个积分 \(\mathbf{u}_{j}=(u_{j},\mathrm{D}u_{j},\ldots,\mathrm{D}^{n-1}u_{j})\) （齐次方程 \(d\mathbf{x}/dt=A(t)\mathbf{x}\) 的积分）线性无关（在空间 \(\mathcal{C}(\mathrm{J};\mathbf{C}^{n})\) ——由 J 到 \(C^{n}\) 的连续映射所组成的空间——中）。显然这个条件是必要的。

反之，若存在 \(n\) 个不全为零的复常数 \(\lambda_j\) ，使得 \(\sum_{j=1}^{n} \lambda_j u_j(t) = 0\) 在 \(J\) 上恒成立，则可推出 \(\sum_{j=1}^{n} \lambda_j D^k u_j(t) = 0\) 在 \(J\) 上对每个整数 \(k\) （满足 \(1 \leqslant k \leqslant n - 1\) ）成立，这就意味着 \(\sum_{j=1}^{n} \lambda_j \mathbf{u}_j(t) = 0\) 在 \(J\) 上成立。

从而 (IV, p.~180, cor. 1)

命题 8. 定义在 J 上的齐次线性方程 (37) 的积分所成之集是域 C 上一个 n 维向量空间。

给定方程 (37) 的任意 n 个积分 \(u_{j}\) ( \(1 \leqslant j \leqslant n\) )，我们把这组积分的朗斯基行列式称为（关于 \(C^{n}\) 的典范基的）相应的 n 个积分 \(u_{j}\) （方程 \(d\mathbf{x}/dt = A(t)\mathbf{x}\) 的积分）的行列式，也就是说，函数

\[
\mathrm{W} (t) = \left| \begin{array}{c c c c} u _ {1} (t) & u _ {2} (t) & \ldots & u _ {n} (t) \\ \mathrm{D} u _ {1} (t) & \mathrm{D} u _ {2} (t) & \ldots & \mathrm{D} u _ {n} (t) \\ \vdots & \vdots & \ldots & \vdots \\ \mathrm{D} ^ {n - 1} u _ {1} (t) & \mathrm{D} ^ {n - 1} u _ {2} (t) & \ldots & \mathrm{D} ^ {n - 1} u _ {n} (t) \end{array} \right|.
\]

为使 n 个积分 \(u_{j}\) 线性无关，必须且只需 \(\mathrm{W}(t) \neq 0\) 在 J 上成立；此外，只需 \(\mathrm{W}(t_{0}) \neq 0\) 在 J 的仅一点 \(t_{0}\) 处成立即可 (IV, p.~184, prop. 4)；进一步，有 (IV, p.~184, prop. 5)

\[
\mathrm{W} (t) = \mathrm{W} (t _ {0}) \exp \left(\int_ {t _ {0}} ^ {t} a _ {1} (s) d s\right).\tag{38}
\]

我们把方程 (35) 的预解 \(C(t, t_{0})\) 与它关于 \(C^{n}\) 的典范基的矩阵等同起来；这个矩阵的各列 \(\mathbf{v}_{j}(t, t_{0})\) ( \(1 \leqslant j \leqslant n\) ) 便是 n 个线性无关的积分

\[
\mathbf {v} _ {j} (t, t _ {0}) = (v _ {j} (t, t _ {0}), \mathrm{D} v _ {j} (t, t _ {0}), \dots , \mathrm{D} ^ {n - 1} v _ {j} (t, t _ {0}))
\]

它们是齐次方程 \(d\mathbf{x}/dt = A(t)\) . x 的积分，对应于方程 (37) 的 n 个线性无关的积分 \(v_{j}(t, t_{0})\) ，这些积分使得

\[
\mathbf {D} ^ {k - 1} v _ {j} (t _ {0}, t _ {0}) = \delta_ {j k}
\]

（克罗内克记号）对 \(1 \leqslant j \leqslant n\) , \(1 \leqslant k \leqslant n\) 成立（约定 \(D^{0}v_{j} = v_{j}\) ）。特别地，由此得出，把常数变易法 (IV, p.~182) 应用于方程 (35)，这里给出 (33) 的一个特积分，它在点 \(t_{0}\) 处连同其前 n - 1 阶导数都等于 0，即函数

\[
w (t) = \int_ {t _ {0}} ^ {t} v _ {n} (t, s) b (s) d s.\tag{39}
\]

在方程 \(D^{n}x = b(t)\) 的特殊情形，公式 (39) 重新给出 \(n^{th}\) 原函数（规则函数 \(b(t)\) 的在点 \(t_{0}\) 处连同其前 n - 1 阶导数都为零者）的表达式，即

\[
w (t) = \int_ {t _ {0}} ^ {t} b (s) \frac {(t - s) ^ {n - 1}}{(n - 1) !} d s
\]

(II, p.~62, 公式 (19))：\(D^{n}x = 0\) 的在点 \(t_{0}\) 处连同其前 n - 2 阶导数都为零、且其 \(n - 1^{th}\) 阶导数在该点等于 1 的积分，实际上就是多项式 \((t - t_{0})^{n-1}/(n - 1)!\) 。

